% =============================================================================
\section{Discussion}\label{sec:discussion}

The renewal L\'evy walk is a convenient idealization, and most of the walkers analyzed here do not follow it. Larvae and termites remember their recent flights, positively and over a few flights, and the memory survives every control we could think of: one field per individual, fields that drift along the track, other ways of cutting tracks into flights, permutation and pipeline nulls, and, for the termites, the wall that they follow. In the larvae it is strongest when brain lobes and sensory neurons have been ablated, which places its origin, together with that of the L\'evy-like statistics~\cite{Sims2019}, below the brain. Part of the dependence between flights may come from switches between behavioral states, which a chain of flights describes as memory; couplings within each state would apportion the two~\cite{Stoye2025}.

Memory changes what a histogram means. A walker that packs long flights into runs has a heavier-tailed histogram than the rule it follows at each turn, by more than half a unit of exponent in the termites. An exponent read from a histogram, often compared with the value $\mu=2$ of optimal search~\cite{Viswanathan1999}, need not be the exponent of the rule that generates the walk. Memory also gives the cutoffs a physical role: the coupling acts through $JL$, the flights condense at the largest scale beyond a self-dual coupling, and over the one or two decades accessible to experiments the walk feels that transition well before reaching it: every termite has crossed its spinodal. Two extensions follow directly: a joint model of the length and the direction of flights, with couplings between speed and turning~\cite{Jerison2020}, and the transport properties of L\'evy walks whose flights are coupled in this way, to which renewal theory does not apply.

In groups, the same formalism separates what the group average mixes: each unit's memory of its own past, drives common to all units, and the influence of the others. Read through the average, own memory appears as a delayed coupling whose strength does not decrease with the size of the group, and a common drive as a coupling at any delay. Read with the self\,+\,social model, the delays found in zebrafish T cells are own memory and a drive shared within each fish, baboons align with where their nearest neighbors, those ahead above all, were heading 10 to 30\,s earlier, and pigeons are coupled to their whole flock from one second to the next. The recipe applies to any attempt to infer delayed interactions from trajectories: model units, not their average; give each unit its own field and memory and each frame its own field; couple units to their neighbors rather than to the group; compare with time-shifted units; and check how a coupling depends on the distance and on the position of the neighbor. Fields fixed in space, such as trails or a homing route, survive both time fields and time shifts, and are the natural next ingredient. Finally, the collective speed has a critical point of its own, at $J_c=12/L^2$, set by the largest susceptibility of a single flight: a delayed positive coupling between scale-free speeds produces collective stop-and-go, a negative one oscillations, and own memory lowers the threshold.
